\section{Money markets and reference rates}

\begin{qfdefinition}[Money-market instruments and reference rates: bills, commercial paper, repo, CDs and bankers' acceptances; fed funds, SOFR, LIBOR's end, EONIA and \texteuro{}STR, SONIA; and the LIBOR--OIS spread as a historical stress gauge]\label{def:money-market-reference-rates}
\qffield{Statement} Money markets are short-term, liquid credit instruments: commercial paper, Treasury bills, agency and municipal paper, bankers' acceptances, time deposits, interbank loans and repurchase agreements \citep[chap.~9, sec.\ ``Funding/Money Markets'', PDF pp.~246--253]{s7-gliner}. In Gliner's account commercial paper is unsecured, ``Most has a maturity of between 90 and 270 days, since it can avoid registering with the Securities and Exchange Commission (SEC) if the security has a maturity of less than 270 days'', and it is cheaper for issuers than bank borrowing (p.~246).

Time deposits (CDs) run from 30 days to five years, those of a year or less counting as money-market instruments (p.~252). Bankers' acceptances are short-term obligations guaranteed by a bank, mainly for trade (pp.~251--252). A repo is ``a sale of securities with a simultaneous agreement by the seller to repurchase them at a later date, which is usually done overnight'', mostly against Treasuries or GSE debt (p.~252).

\emph{Fed funds} are unsecured overnight loans of reserve balances, and the effective rate serves as the reference for OIS (p.~257). \emph{LIBOR}, in the BBA's definition, was the rate at which a panel bank could borrow ``by asking for and then accepting interbank offers in reason [sic] market size, just prior to 11:00 A.M. London time'', computed as a trimmed mean of 16 banks' submissions (p.~254). The LIBOR--OIS spread compares unsecured term bank borrowing with the compounded overnight rate; it ``skyrocketed'' in 2008, rose with European banks' dollar funding needs in the euro crisis, and ``does not automatically indicate a crisis'' but is ``a useful barometer for credit risk'' (pp.~261--262).

The SEC staff's interpretations add that commercial paper ``payable on demand but in any event no later than 9 months from issuance'' meets the maturity requirement, and that tracing proceeds to actual current transactions is not necessary where they are commingled with general funds and an equal amount is applied to current transactions (Corporation Finance interpretations 219.01 and 219.02, both dated 26 November 2008) \citep[Securities Act Sections, 219.01--219.02, page updated 23 June 2026]{sec-cdi-securities-act}.

The Richmond Fed's money-market handbook (chapter last updated 1993) lists further historical conditions: notes ``of a type not ordinarily purchased by the general public'' (in practice minimum denominations usually \$100,000), current-transaction use, and, under SEC Release 33-4412 (1961), ``prime quality'' and discountability at a Federal Reserve Bank, the discounting requirement having been ``dropped in 1980''; it puts typical maturities at 5 to 45 days, averaging 30--35, against Gliner's 90 to 270, and notes that rollover is allowed if not automatic \citep[chap.~9, pp.~105--107 and n.~4]{richmondfed-cp-1993}. Private placements (Section 4(a)(2)) and bank-guaranteed paper (Section 3(a)(2)) are other exemptions.

\qffield{Notation} SOFR the Secured Overnight Financing Rate; EFFR the effective federal funds rate; \texteuro{}STR the euro short-term rate; SONIA the Sterling Overnight Index Average.

\qffield{Conditions} \emph{Current benchmarks.} SOFR measures the cost of overnight borrowing collateralised by Treasuries, and the EFFR is a volume-weighted median of overnight federal funds transactions, both published by the New York Fed \citep[SOFR description; EFFR and OBFR methodology]{exp-sofr-data}. The U.S.\ dollar LIBOR panel ended after 30 June 2023 \citep[web page]{arrc-sofr-transition}, with synthetic settings to 30 September 2024 \citep[sec.~2, n.~11]{cbot-23-214}. EONIA was discontinued on 3 January 2022 \citep[web page]{emmi-eonia}; the ECB first published \texteuro{}STR on 2 October 2019 \citep[web page]{ecb-estr}. A LIBOR--OIS spread therefore exists only historically.

\emph{One documented replacement}. The Office of Financial Research's Financial Stress Index replaced all seven variables in its funding category from 1 January 2022 (published 27 June 2023): 3-month U.S.\ LIBOR--OIS by ``3-Month Commercial Paper - SOFR OIS'', the TED spread by ``3-Month Commercial Paper - 3-Month Treasury'', the 2-year LIBOR swap spread by the 2-year SOFR swap spread, 3-month Euribor--EONIA OIS by Euribor--\texteuro{}STR OIS, and the yen and euro cross-currency spreads by SOFR, \texteuro{}STR and TONAR versions, with the commercial-paper rate the Federal Reserve Board's 3-month AA financial rate; the index's methodology was unchanged \citep[pp.~1--2, Tables~1--2]{ofr-fsi-factsheet-2023}. In the index's daily data, the funding category's contribution was 9.58 on 10 October 2008 and 0.96 on 30 March 2020 (both with the old variables), and $-0.16$ on 13 October 2026, below its long-run average of zero \citep[\texttt{fsi.csv}]{ofr-fsi-data}.

That is one official monitor's choice, not a survey of what market studies use; the Federal Reserve's May 2026 \emph{Financial Stability Report} judges funding risk from quantities (runnable liabilities, uninsured deposits, liquid assets against short-term debt), not from a spread \citep[sec.~4]{fed-fsr-2026-05}.

\qffield{Practice} Treasury sells bills ``for terms ranging from four weeks to 52 weeks'', auctioning 4-, 6-, 8-, 13-, 17- and 26-week bills weekly and 52-week bills every four weeks, plus cash management bills on no regular schedule \citep[Treasury Bills page]{treasurydirect-bills} (Gliner gives the history of bill auctions, p.~247). That money-market stress often precedes broader financial stress is a plausible but unsourced generalisation.

\begin{itemize}
\item \emph{Commercial-paper maturities.} The BIS primer gives the legal limit as ``270 days in the United States and one year in Europe and Japan'', but ``In practice, the maturity of CP tends to be much shorter. In the US, the average maturity is around 30 days, while in Europe it is approximately four months'' \citep[sec.\ ``Main funding instruments and their characteristics'']{gap-money-cp-cd-primer}. Gliner's ``Most has a maturity of between 90 and 270 days'' is therefore not a description of the U.S.\ market as the BIS reports it; the 270-day figure is the registration exemption, not typical practice.
\item \emph{How SOFR is calculated and revised.} SOFR covers the trades in the BGCR (tri-party and GCF) plus cleared bilateral (DVP) repo; to limit the influence of specials, ``20 percent of the lowest-rate transaction volume from the DVP segment is removed each day''; the rate is the volume-weighted median, rounded to the basis point, published at about 8:00 a.m.\ ET with the 1st, 25th, 75th and 99th volume-weighted percentiles and the volume; a revision can follow at about 2:30 p.m.\ ET the same day, and ``Rate revisions will only occur if the change in the rate exceeds one basis point and only on the same day as initial publication'' \citep[sections on SOFR, calculation methodology and publication and revisions]{exp-sofr-data}. A backtest that uses SOFR as of the morning should record whether the rate was later revised that day.
\item \emph{Why overnight RFRs.} Schrimpf and Sushko trace the switch to the collapse of unsecured interbank term lending: few transactions underpinned LIBOR even before 2008, and interbank trading ``plummeted'' after it. RFR-linked OIS give a risk-free term structure for discounting, but banks lose a benchmark for their marginal term funding cost. Backward-looking (compounded in arrears) term rates are known only at the end of the period and lag rising rates, while forward-looking term rates from futures or OIS are known in advance \citep[BIS Quarterly Review, March 2019, sections on the case for reform and on term rates]{rt-schrimpf-sushko-rfr-2019}.
\item \emph{A liquidity premium inside repo rates.} Copeland and Engbretson measure the premium dealers charge for holding liquid buffers: at 2024's average rates and reserves, ``a 100-basis-point increase in the interest rate on reserve balances results in a 0.9 basis point increase in the liquidity risk premium'', about 10 percent of the dealer spread, the same effect as a \$429 billion fall in reserves, and the relation is nonlinear \citep[Staff Report~1189, March 2026, abstract]{latest-copeland-repo-premium-2026}.
\end{itemize}

\qffield{Pitfalls} LIBOR did not end simply because of the manipulation scandals: the cases undermined confidence, but the cessation followed the lack of underlying unsecured transactions and the FCA's 2021 announcements (\hyperref[meth:policy-path-ois-futures]{policy path}); U.S.\ dollar panel settings continued to 30 June 2023. \texteuro{}STR began in 2019 and EONIA ended on 3 January 2022. SONIA was reformed in 2018 to cover more transactions. Gliner describes a pre-2013 LIBOR calculation by the BBA; administration passed to ICE Benchmark Administration later.

\qffield{See also} \hyperref[def:interest-rate-measures]{interest-rate measures}; \hyperref[meth:policy-path-ois-futures]{policy path from OIS and futures}; \hyperref[mech:securities-lending-short-sale]{securities lending}; \hyperref[def:fixed-income-universe]{fixed-income universe}.

\qffield{Further reading} On funding liquidity, \citep{exp-brunnermeier-pedersen-liquidity}.

\end{qfdefinition}

\begin{qfevidence}[Runs on prime money market funds in March 2020: institutional investors ran on fundamentals, retail investors followed within-family outflows]\label{ev:money-fund-runs}
\qffield{Statement} At the height of the March 2020 run, cumulative redemptions from U.S.-dollar prime money market funds, onshore and offshore, were 22 percent of the industry's total net assets at end-2019: 32 percent for institutional and 11 percent for retail investors \citep[Staff Report~956, December 2020, revised April 2024, sec.~1, PDF p.~3]{gap-money-prime-fund-runs}. Cipriani and La Spada find that ``sophisticated institutional investors ran preemptively based on fundamentals'' while ``unsophisticated retail investors ran based on herd-like informational spillovers'', leaving funds of families with large institutional outflows; institutional investors looked up fund portfolios more (website traffic), and both types ran more when switching to a safer fund was cheaper, ``suggesting that safe-haven availability exacerbates runs'' (abstract, PDF p.~2).

\qffield{Conditions} One episode and market; sophistication is split by share class, and website traffic proxies for information acquisition.

\qffield{Practice} For a run scenario, model institutional and retail outflows separately, and allow outflows to depend on a cheap switch into government funds. The 2023 SEC reforms changed liquidity fees and gates (\nolinkurl{gap-money-sec-2023-reform}).

\qffield{Pitfalls} Figures are from the working paper's abstract and introduction.

\qffield{See also} \hyperref[def:money-market-reference-rates]{money-market instruments and reference rates}; \hyperref[ev:stablecoin-runs]{stablecoin runs and flights to safety}.

\qffield{Further reading} \nolinkurl{gap-money-sec-2023-reform}; \nolinkurl{gap-money-nmfp-data}.

\end{qfevidence}
